
\documentclass{article}
\usepackage{colortbl}
\usepackage{makecell}
\usepackage{multirow}
\usepackage{supertabular}

\begin{document}

\newcounter{utterance}

\centering \large Interaction Transcript for game `referencegame', experiment `line\_grids\_rows', episode 13 with reflexion{-}Qwen3.5{-}4B.
\vspace{24pt}

{ \footnotesize  \setcounter{utterance}{1}
\setlength{\tabcolsep}{0pt}
\begin{supertabular}{c@{$\;$}|p{.15\linewidth}@{}p{.15\linewidth}p{.15\linewidth}p{.15\linewidth}p{.15\linewidth}p{.15\linewidth}}
    \# & $\;$A & \multicolumn{4}{c}{Game Master} & $\;\:$B\\
    \hline

    \theutterance \stepcounter{utterance}  
    & & \multicolumn{4}{p{0.6\linewidth}}{
        \cellcolor[rgb]{0.9,0.9,0.9}{
            \makecell[{{p{\linewidth}}}]{
                \texttt{\tiny{[P1$\langle$GM]}}
                \texttt{You have attempted this same task before and did not succeed. These are the plans you wrote after those attempts.} \\
\texttt{Trial \#1: New plan: I failed because my description was inaccurate regarding the number of filled cells and their positions. The target grid has a vertical column of **five** filled cells in the fourth column, not three. I should have carefully counted the rows and columns before generating the expression. Next time, I will describe the target grid as "The grid with a vertical column of five filled cells in the fourth column" to precisely distinguish it from the distractors which have fewer filled cells or different column positions.} \\
\\ 
\texttt{Trial \#2: New plan: I failed because my description was inaccurate regarding the number of filled cells and their positions. The target grid has a vertical column of **five** filled cells in the fourth column, not three. I should have carefully counted the rows and columns before generating the expression. Next time, I will describe the target grid as "The grid with a vertical column of five filled cells in the fourth column" to precisely distinguish it from the distractors which have fewer filled cells or different column positions.} \\
\\ 
\\ 
\texttt{You are given three grids, where each of them is 5 by 5 in size.} \\
\texttt{Grids have empty cells marked with "▢" and filled cells marked with "X".} \\
\texttt{Your task is to generate a referring expression that best describes the target grid while distinguishing it from the two other distractor grids.} \\
\texttt{The first grid is the target grid, and the following two grids are the distractors.} \\
\\ 
\texttt{Target grid:} \\
\\ 
\texttt{X X □ X X} \\
\texttt{X X □ X X} \\
\texttt{X X □ X X} \\
\texttt{X X □ X X} \\
\texttt{X X □ X X} \\
\\ 
\texttt{Distractor grid 1:} \\
\\ 
\texttt{X □ □ X X} \\
\texttt{X □ □ X X} \\
\texttt{X □ □ X X} \\
\texttt{X □ □ X X} \\
\texttt{X □ □ X X} \\
\\ 
\texttt{Distractor grid 2:} \\
\\ 
\texttt{□ □ □ X X} \\
\texttt{□ □ □ X X} \\
\texttt{□ □ □ X X} \\
\texttt{□ □ □ X X} \\
\texttt{□ □ □ X X} \\
\\ 
\texttt{Instruction: Describe the target grid.} \\
\texttt{Generate the referring expression starting with the tag "Expression: " for the given target grid. Omit any other text.} \\
            }
        }
    }
    & \\ \\

    \theutterance \stepcounter{utterance}  
    & \multicolumn{4}{p{0.6\linewidth}}{
        \cellcolor[rgb]{0.9,0.9,0.9}{
            \makecell[{{p{\linewidth}}}]{
                \texttt{\tiny{[P1$\rangle$GM]}}
                \texttt{Expression: The grid with a vertical column of five filled cells in the fourth column.} \\
            }
        }
    }
    & & \\ \\

    \theutterance \stepcounter{utterance}  
    & & & \multicolumn{2}{p{0.3\linewidth}}{
        \cellcolor[rgb]{0.9,0.9,0.9}{
            \makecell[{{p{\linewidth}}}]{
                \texttt{\tiny{[GM$|$GM]}}
                \texttt{Expression: The grid with a vertical column of five filled cells in the fourth column.} \\
            }
        }
    }
    & & \\ \\

    \theutterance \stepcounter{utterance}  
    & & \multicolumn{4}{p{0.6\linewidth}}{
        \cellcolor[rgb]{0.9,0.9,0.9}{
            \makecell[{{p{\linewidth}}}]{
                \texttt{\tiny{[GM$\rangle$P2]}}
                \texttt{You have attempted this same task before and did not succeed. These are the plans you wrote after those attempts.} \\
\texttt{Trial \#1: New plan: I misidentified the target grid by focusing on the wrong column. The expression specifies the "fourth column," but I likely matched the pattern in the second column of the First grid instead. Next time, I will explicitly count the columns from left to right (1, 2, 3, 4, 5) and verify that the filled cells ("X") appear exactly in the fourth column before selecting the grid. I will also ensure I count the rows correctly to confirm the vertical sequence of three filled cells in the first two rows and one in the third.} \\
\\ 
\texttt{Trial \#2: New plan: I misidentified the target grid by confusing the column indices. The expression specifies the "fourth column," but I likely selected the First grid because its filled cells appeared in the second column, which I incorrectly mapped to the fourth. Next time, I will explicitly count the columns from left to right (1, 2, 3, 4, 5) and verify that the filled cells ("X") appear exactly in the fourth column before selecting the grid. I will also ensure I count the rows correctly to confirm the vertical sequence of five filled cells in the first two rows and one in the third.} \\
\\ 
\\ 
\texttt{You are given three grids, where each of them is 5 by 5 in size.} \\
\texttt{Grids have empty cells marked with "▢" and filled cells marked with "X".} \\
\texttt{You are also given a referring expression that describes one of the given grids.} \\
\texttt{Your task is to select the grid that matches the given referring expression.} \\
\texttt{Generate only the number (in text) of the grid that matches the given expression by selecting first, second, or third.} \\
\\ 
\texttt{First:} \\
\\ 
\texttt{X □ □ X X} \\
\texttt{X □ □ X X} \\
\texttt{X □ □ X X} \\
\texttt{X □ □ X X} \\
\texttt{X □ □ X X} \\
\\ 
\texttt{Second:} \\
\\ 
\texttt{X X □ X X} \\
\texttt{X X □ X X} \\
\texttt{X X □ X X} \\
\texttt{X X □ X X} \\
\texttt{X X □ X X} \\
\\ 
\texttt{Third:} \\
\\ 
\texttt{□ □ □ X X} \\
\texttt{□ □ □ X X} \\
\texttt{□ □ □ X X} \\
\texttt{□ □ □ X X} \\
\texttt{□ □ □ X X} \\
\\ 
\texttt{Expression: The grid with a vertical column of five filled cells in the fourth column.} \\
\texttt{Question: Which grid does the expression refer to?} \\
\texttt{Start with the tag "Answer: ", followed by your selection. Omit any other text.} \\
            }
        }
    }
    & \\ \\

    \theutterance \stepcounter{utterance}  
    & & & \multicolumn{4}{p{0.6\linewidth}}{
        \cellcolor[rgb]{0.9,0.9,0.9}{
            \makecell[{{p{\linewidth}}}]{
                \texttt{\tiny{[GM$\langle$P2]}}
                \texttt{Answer: Third} \\
            }
        }
    }
     \\ \\

    \theutterance \stepcounter{utterance}  
    & & & \multicolumn{2}{p{0.3\linewidth}}{
        \cellcolor[rgb]{0.9,0.9,0.9}{
            \makecell[{{p{\linewidth}}}]{
                \texttt{\tiny{[GM$|$GM]}}
                \texttt{third} \\
            }
        }
    }
    & & \\ \\

    \theutterance \stepcounter{utterance}  
    & \multicolumn{4}{p{0.6\linewidth}}{
        \cellcolor[rgb]{0.95,0.95,0.95}{
            \makecell[{{p{\linewidth}}}]{
                \texttt{\tiny{[P2$\rightarrow$P1]}}
                \texttt{Post{-}episode Reflexion (not sent to GM):} \\
\texttt{New plan: I failed because my description was inaccurate regarding the number of filled cells and their positions. The target grid has a vertical column of **five** filled cells in the fourth column, not three. I should have carefully counted the rows and columns before generating the expression. Next time, I will describe the target grid as "The grid with a vertical column of five filled cells in the fourth column" to precisely distinguish it from the distractors which have fewer filled cells or different column positions.} \\
            }
        }
    }
    & & \\ \\

    \theutterance \stepcounter{utterance}  
    & & & \multicolumn{4}{p{0.6\linewidth}}{
        \cellcolor[rgb]{0.95,0.95,0.95}{
            \makecell[{{p{\linewidth}}}]{
                \texttt{\tiny{[P1$\rightarrow$P2]}}
                \texttt{Post{-}episode Reflexion (not sent to GM):} \\
\texttt{New plan: I misidentified the target grid by confusing the column indices and the number of filled cells. The expression specifies a "vertical column of five filled cells," but I likely matched the First grid's pattern of four filled cells instead of the Third grid's pattern of five. Next time, I will explicitly count the columns from left to right (1, 2, 3, 4, 5) and verify that the filled cells ("X") appear exactly in the fourth column. Crucially, I will count the total number of filled cells in that column to ensure it is exactly five, matching the "vertical column of five" description, before selecting the grid.} \\
            }
        }
    }
     \\ \\

\end{supertabular}
}

\end{document}
